% Lunette astronomique afocale : marche d'un faisceau issu d'un point à l'infini (L'essentiel, chapitre 24)
% objectif : f'1 = 6 ; oculaire : f'2 = 1,5 ; tan(theta) = 0,1 ; tan(theta') = 0,4
\begin{tikzpicture}[x=1cm,y=1cm]
\draw[otMuted,line width=0.4pt] (-2.4,0) -- (10.4,0);
\draw[otNavy,line width=1pt,{Stealth[length=2mm]}-{Stealth[length=2mm]}] (0,-1.6) -- (0,1.6) node[pcLegende,anchor=south,text=otNavy] {objectif $L_1$};
\draw[otNavy,line width=1pt,{Stealth[length=2mm]}-{Stealth[length=2mm]}] (7.5,-1.3) -- (7.5,1.3) node[pcLegende,anchor=south,text=otNavy] {oculaire $L_2$};
\fill[otInk] (0,0) circle (0.04) node[pcLegende,anchor=north east] {$O_1$};
\fill[otInk] (7.5,0) circle (0.04) node[pcLegende,anchor=north west] {$O_2$};
\fill[otInk] (6,0) circle (0.04) node[pcLegende,anchor=south east] {$F'_1=F_2$};
\fill[otInk] (9,0) circle (0.04) node[pcLegende,anchor=north] {$F'_2$};
\draw[otMuted,line width=0.4pt,dash pattern=on 2pt off 1.5pt] (6,-1.2) -- (6,0.35);
% faisceau incident (pente -0,1) et convergence en B1
\foreach \y in {1.2,0,-0.6}{
  \draw[otRed!80!black,line width=0.7pt] (-2.2,{\y+0.22}) -- (0,\y) -- (6,-0.6);
  \draw[otRed!80!black,line width=0.7pt,-{Stealth[length=1.5mm]}] (-2.2,{\y+0.22}) -- (-1.1,{\y+0.11});}
% sortie de l'oculaire : rayons parallèles de pente 0,4
\foreach \yd/\pente in {-1.05/-0.3,-0.75/-0.1,-0.6/0}{
  \draw[otRed!80!black,line width=0.7pt] (6,-0.6) -- (7.5,\yd);
  \draw[otRed!80!black,line width=0.7pt,-{Stealth[length=1.5mm]}] (7.5,\yd) -- (10.2,{\yd+1.08});}
\draw[otMuted,line width=0.4pt,dash pattern=on 1.5pt off 1pt] (6,-0.6) -- (7.5,0);
\fill[otRed!80!black] (6,-0.6) circle (0.05);
\node[pcLegende,anchor=north,text=otRed!80!black] at (5.9,-0.65) {$B_1$};
\draw[otMuted,line width=0.4pt] (-1.0,0) arc[start angle=180,end angle=174.3,radius=1.0];
\node[pcLegende,anchor=south east] at (-1.0,0.02) {$\theta$};
\draw[otMuted,line width=0.4pt] (9.3,0) arc[start angle=0,end angle=21.8,radius=0.9];
\node[pcLegende,anchor=west] at (9.32,0.2) {$\theta'$};
\draw[otMuted,line width=0.4pt,dash pattern=on 2pt off 1.5pt] (8.4,0) -- (9.6,0.48);
\draw[otMuted,line width=0.4pt,{Stealth[length=1.2mm]}-{Stealth[length=1.2mm]}] (0,-1.85) -- (6,-1.85) node[pos=0.5,pcLegende,anchor=north] {$f'_1$};
\draw[otMuted,line width=0.4pt,{Stealth[length=1.2mm]}-{Stealth[length=1.2mm]}] (6,-1.85) -- (7.5,-1.85) node[pos=0.5,pcLegende,anchor=north] {$f'_2$};
\end{tikzpicture}
\par\smallskip\legende{L'objectif forme l'image intermédiaire $B_1$ dans son plan focal image, confondu avec le plan focal objet de l'oculaire : le faisceau émergent est parallèle, incliné de $\theta'$, parallèle au rayon $B_1O_2$.}
